% Avsnitt 7.1-7.5, ur lectures/L07/appendix/a_timers.md (A.1-A.5).

\section{Timers: en räknare med ett målvärde}
\label{c7:sec:concept}

En \textbf{timer} svarar på en annan fråga än räknaren från \secref{c6:sec:counter}. En räknare
svarar på ”vad är räkningen just nu?”; en timer svarar på ”har en viss tid gått?”

Den gör det med tre delar: en intern räknare som är osynlig utifrån, ett målvärde
\code{TICK_COUNT} som är fast per instans, och en jämförelse som när den är sann nollställer
räknaren och pulsar utgångsflaggan \code{timeout} under en klockcykel.

Eftersom systemklockan tickar med en känd, fast frekvens är att räkna tick samma sak som att mäta
förfluten tid. På DE0-CV är det 50\,000\,000 tick per sekund:

\[ \mathit{TICK\_COUNT} = \mathit{sekunder} \times 50\,000\,000 \]

\begin{itemize}
  \item \code{TICK_COUNT = 50_000_000}: timeout en gång per sekund.
  \item \code{TICK_COUNT = 25_000_000}: två gånger per sekund.
  \item \code{TICK_COUNT = 5_000_000}: tio gånger per sekund.
\end{itemize}

Strikt taget löper räknaren \code{0} till och med \code{TICK_COUNT}, så en full period är
\code{TICK_COUNT + 1} cykler. Det där extra ticket är mellan 0,02 och 0,2~ppm vid de här värdena,
långt under kortoscillatorns egen tolerans, så de runda talen ovan används för tydlighetens skull;
för ett exakt cykelantal, använd
$\mathit{TICK\_COUNT} = \mathit{sekunder} \times 50\,000\,000 - 1$.

Det här är motsatsen till en blockerande \code{_delay_ms()}. En blockerande fördröjning stoppar hela
programmet för att vänta; en timer går permanent parallellt med allt annat och höjer helt enkelt en
flagga, och ingenting i designen blockeras någonsin i väntan på den.

% ------------------------------------------------------------------------------------------
\section{Att bygga en timer för hand i CircuitVerse}
\label{c7:sec:byhand}

Innan någon VHDL skrivs: bygg en timer som ett grindnät, precis som du gjorde för \chapref{c1:ch}
och \ref{c2:ch}:s kombinatoriska nät och \chapref{c3:ch}:s vippor. En timer är liten nog att rita i
sin helhet, och det är att rita den som gör att \secref{c7:sec:module}:s VHDL läses som en
beskrivning av något du redan förstår snarare än ett nytt idiom att lära utantill.

Bygg den i en storlek du kan följa: en \textbf{4-bitars räknare} med \code{TICK_COUNT = 10}, vilket
är \code{1010}. Den riktiga designen räknar till 50\,000\,000; ingenting i strukturen ändras, bara
bredden.

\textbf{Utgå från räknaren du redan byggt} i \secref{c6:sec:counter}: ett 4-bitars register
\code{Q0} till \code{Q3} som delar en klocka, en adderare och en konstant \code{1} med summan
återkopplad till \code{D}-ingångarna, och ingenting mer. Om du inte sparade den är det tre element
och ungefär fem minuter att rita om. Bekräfta att den fortfarande räknar och slår över innan du
lägger till något, så att du vet att räknaren under är sund när timern senare krånglar.

Allt här lägger \textbf{två saker ovanpå det}: en jämförelse som säger ”räkningen är framme”, och en
nollställning som får den att upprepa. Det är hela skillnaden mellan en räknare och en timer.

\paragraph{Lägg till jämförelsen}
Du vill ha \code{timeout} hög exakt när de fyra bitarna stavar \code{1010}, och det rena sättet att
bygga ”är de här två talen lika?” är en \textbf{XNOR-grind per bit}. En XNOR ger \code{1} när dess
ingångar är \emph{lika}, vilket är \secref{c1:sec:gates}:s XNOR-kolumn läst som ett likhetstest.
Koppla in fyra, var och en som jämför en räknarutgång mot motsvarande målbit (\code{Q3} mot
\code{1}, \code{Q2} mot \code{0}, \code{Q1} mot \code{1}, \code{Q0} mot \code{0}), och driv
målsidan från en konstant i CircuitVerse snarare än från en ledning du senare kan råka ta för en
signal. Kombinera sedan de fyra utgångarna med AND; den enda AND-grinden är \code{timeout}. Med
$\odot$ för XNOR, den enda operator \secref{c1:sec:boolean} inte gav någon symbol:

\[ \mathit{timeout} = (Q_3 \odot 1) \cdot (Q_2 \odot 0) \cdot (Q_1 \odot 1) \cdot (Q_0 \odot 0) \]

Alla fyra bitar måste stämma samtidigt, vilket är vad ”räknaren har nått \code{1010}” betyder. Det
här är en \textbf{allmän likhetsjämförare}, värd att känna igen som ett eget byggblock: samma fyra
XNOR plus en AND jämför mot vilket 4-bitars värde som helst, och bara konstanterna ändras.

\paragraph{Lägg till nollställningen}
Så att timern startar om i stället för att fortsätta: koppla tillbaka \code{timeout} för att tvinga
räknaren till \code{0000} vid nästa flank, genom att grinda varje vippas \code{D}-ingång så att
\code{timeout = 1} går före adderarens utgång. Utan den fortsätter räknaren till \code{1111} och
slår över av sig själv, så \code{timeout} slår till var sextonde flank i stället för var elfte. Båda
är periodiska; bara den ena har den period du bad om.

\paragraph{Bekräfta sedan, en klockflank i taget}
\begin{itemize}
  \item Räknaren går från \code{0000} till och med \code{1010}, och \code{timeout} är låg för varje
    räknarvärde utom \code{1010}. Håll ögonen på de fyra XNOR-utgångarna medan den räknar och lägg
    märke till hur ofta tre är höga samtidigt. Tre bitar som stämmer är ingen träff; det är
    AND-grinden som insisterar på alla fyra.
  \item Vid flanken efter att \code{timeout} slagit till är räknaren tillbaka på \code{0000} och
    \code{timeout} är låg, så den är hög under exakt en klockcykel, aldrig två.
  \item Hela cykeln är alltså elva flanker, räknarvärdena \code{0} till och med \code{10}, upprepade
    i all oändlighet. Det är \secref{c7:sec:concept}:s ettfel \code{TICK_COUNT + 1}, och det här är
    bokens enklaste ställe att se det på, eftersom du kan räkna flankerna med ögat.
\end{itemize}

\lecfig[0.5]{lectures/L07/appendix/images/timer_circuit.png}{Jämförelselogiken: fyra XNOR mot
\code{1010}, ANDade till \code{timeout}.}{c7:fig:timercircuit}

Ritningen är beskuren till den del som det här avsnittet lägger till. Den märker räknarens bitar
\code{counter[3]} ned till \code{counter[0]} där texten ovan kallar dem \code{Q3} till \code{Q0},
och uppräkningen som matar dem är ritad som carry-logik på grindnivå snarare än det enda
adderarblock \secref{c6:sec:counter} använde, eftersom det är vad CircuitVerse ger dig när du bygger
den en bit i taget. De fyra XNOR-grindarna, deras konstanter \code{1}, \code{0}, \code{1},
\code{0}, och den AND som kombinerar dem är allt som är nytt här.

Två saker är värda att lägga märke till före VHDL:en, eftersom båda är vad språket köper dig:
\begin{itemize}
  \item \textbf{Jämföraren är fast kopplad till ett målvärde.} Att ändra \code{TICK_COUNT} från
    \code{10} till \code{12} innebär att gå tillbaka in i ritningen och vända på två konstanter. Det
    är en liten ändring, men det är en ändring i \emph{kretsen}, och varje instans behöver sin egen
    kopia. I \secref{c7:sec:module} blir målvärdet en generic.
  \item \textbf{Bredden är fast kopplad också.} Att räkna till 50\,000\,000 kräver 26 vippor och en
    26-bitars jämförare, 26 XNOR över 52 ingångar: samma krets, en orimlig ritning. I
    \secref{c7:sec:module} är det samma sex rader VHDL hur som helst.
\end{itemize}

Det är avvägningen den här boken gör hela tiden: rita kretsen en gång, i en storlek du kan se, för
att veta vad verktyget bygger åt dig, och låt sedan verktyget bygga den i en storlek du inte hade
kunnat rita.

% ------------------------------------------------------------------------------------------
\section{Modulen \code{timer} i VHDL}
\label{c7:sec:module}

Samma krets som en återanvändbar modul. Den byggs live under passet, och
Övning~\ref{c7:ex:timer} ber dig skriva den själv efteråt, så den är tryckt här i sin helhet
snarare än utdelad som en fil:

\begin{vhdlcode}
entity timer is
    generic(TICK_COUNT: natural := 50_000_000);
    port(clock, reset_s2_n, enable: in std_logic;
         timeout                  : out std_logic);
end entity;

architecture behaviour of timer is
-- Internal tick counter.
signal counter: natural range 0 to TICK_COUNT;
begin
    process(clock, reset_s2_n) is
    begin
        if (reset_s2_n = '0') then
            timeout <= '0';
            counter <= 0;
        elsif (rising_edge(clock)) then
            timeout <= '0';
            if (enable = '1') then
                if (TICK_COUNT > counter) then
                    counter <= counter + 1;
                else
                    timeout <= '1';
                    counter <= 0;
                end if;
            end if;
        end if;
    end process;
end architecture;
\end{vhdlcode}

\begin{itemize}
  \item \code{TICK_COUNT} är en \textbf{generic}, inte en konstant inbakad i arkitekturen, så samma
    entitet betjänar vilken frekvens som helst genom att instansieras med en annan
    \code{generic map}. Den interna räknarens intervall är skrivet i termer av den, så registret
    dimensioneras efter målvärdet snarare än efter vad ett obegränsat \code{natural} skulle ge,
    samma poäng som \secref{c6:sec:counter} gör om \code{natural range 0 to 15}. En generic kan
    dimensionera en signal just för att den är fastlåst innan syntesen körs.
  \item \code{reset_s2_n} är den redan \emph{synkroniserade} reset-signalen, aldrig den råa
    asynkrona \code{reset_n}. En timer rör aldrig en rå asynkron ingång direkt.
  \item \code{enable} \textbf{fryser} timern snarare än nollställer den. Medan den är \code{'0'}
    håller \code{counter} sitt värde och \code{timeout} förblir låg, och en avstängd timer
    fortsätter där den slutade snarare än startar om. \Chapref{c8:ch}:s \code{fsm_led} hänger på
    det.
  \item \code{timeout} är en encykelspuls. Den \code{else}-gren som sätter den nås vid bara en
    flank: den flank där \code{counter} redan står på \code{TICK_COUNT} och nollställs tillbaka till
    \code{0}. Vid den följande flanken körs den villkorslösa \code{timeout <= '0';} utan något efter
    sig som kan skriva över, så pulsen är exakt en cykel bred.
\end{itemize}

Jämför det här med \secref{c7:sec:byhand}:s grindnät. Logiken är densamma, och
\secref{c7:sec:byhand} namngav redan vad genericen köper. En skillnad namngav den inte:
\secref{c7:sec:byhand}:s \code{timeout} är en kombinatorisk avkodning, hög under den cykel då
räkningen visar \code{1010}, medan den här är registrerad, hög en cykel senare, under den cykel då
\code{counter} visar \code{0}. Samma period, annan fas, och annat glitchbeteende, vilket är den
distinktion som \chapref{c8:ch}:s diskussion om Moore kontra Mealy vilar på.

% ------------------------------------------------------------------------------------------
\section{Att komponera en komplett timerkrets}
\label{c7:sec:compose}

En timer ensam pulsar bara \code{timeout}; något måste reagera på den pulsen, och en knapp måste
styra timern. Med \chapref{c4:ch}:s dubbelvippsynkroniserare återanvänd ser en komplett knappstyrd,
timerdriven lysdiodskrets ut så här:

\lecfig[1.0]{lectures/L07/appendix/images/led_toggle_timer.png}{Hela kretsen på en gång:
synkroniserad knapp, timer och växlad lysdiod. Detaljerna är läsbara i de två närbilderna
nedan.}{c7:fig:whole}

Läs den som en planritning snarare än som ett kopplingsschema. Det är en bred krets, så i
sidstorlek blir dess etiketter små; vad den är till för är att visa hur få delar det finns och
ungefär var var och en sitter. De två närbilderna nedan är där signalnamnen går att läsa, och
tillsammans täcker de allt i den.

Från vänster till höger:
\begin{itemize}
  \item \textbf{Synkronisering och flankdetektering.} \code{button_n} passerar två vippor för
    metastabilitetsskydd, och en tredje låter sedan kretsen jämföra ”nu” mot ”föregående” och
    detektera en fallande flank, vilket ger en \code{button_edge_s2}-puls på en cykel.
  \item \textbf{Timern} (\secref{c7:sec:module}), aktiverad eller avstängd av
    \code{timer_enabled}, som i sin tur växlas av \code{button_edge_s2}.
  \item \textbf{Att växla lysdioden.} Timerns \code{timeout}-puls vänder lysdiodens läge varje gång
    den slår till, och lysdioden tvingas släckt så snart timern är avstängd eller kretsen resettas.
\end{itemize}

\lecfig[0.98]{lectures/L07/appendix/images/button_sync_and_flank_detection.png}{Synkroniserare och
detektor för fallande flank.}{c7:fig:buttonsync}

\lecfig[0.6]{lectures/L07/appendix/images/led_toggle.png}{Växlingslogiken för lysdioden, driven av
timerns \code{timeout}-puls.}{c7:fig:ledtoggle}

Två mönster spelar roll här, oberoende av den här kretsen:
\begin{itemize}
  \item \textbf{Varje asynkron ingång synkroniseras innan den rör någon annan logik}, och varje
    enable- eller växlingssignal som härleds ur en knapp är en synkron signal från den punkten och
    framåt. Du återanvänder det oförändrat i \secref{c7:sec:walking} och igen för \chapref{c8:ch}:s
    tillståndsmaskiner.
  \item \textbf{En klocka, och allt annat är enable.} \code{timeout} taktar lysdioden genom att
    vara en \emph{enable} på en cykel för en vippa som går på systemklockan, inte genom att kopplas
    in i en klockingång. Det är frestande att dra en räknar- eller timerutgång in i en klockport för
    att ”sakta ned”; låt bli. En design med en klocka är en design vars tajming verktyget kan
    analysera, och det är det enskilt vanligaste misstaget en programmerare gör när en
    baudgenerator ska skrivas.
\end{itemize}

% ------------------------------------------------------------------------------------------
\section{\code{walking_led}: ett skiftregister med vandrande lysdiod}
\label{c7:sec:walking}

\code{walking_led} är den första designen som komponerar tre kapitels byggblock till en krets som
går att demonstrera i hårdvara: en enda tänd lysdiod som vandrar längs en rad, en position per
timertick, startad och stoppad med en knapp. Den byggs live under passet, och
Övning~\ref{c7:ex:walking} ber dig skriva den själv efteråt, så den är tryckt här i sin helhet
snarare än utdelad som en fil.

Den skriver nästan ingen ny logik, utan återanvänder \code{reset_sync} och \code{button_sync}
oförändrade (\secref{c4:sec:resetsync} och \ref{c4:sec:buttonsync}) för att synkronisera och
flankdetektera knappen, \code{timer} (\secref{c7:sec:module}) för att takta hastigheten, och
skiftregistret från \secref{c6:sec:shift}--\ref{c6:sec:piso} för själva vandringen. Det är det de
fyra senaste kapitlen har samlat på sig: en liten uppsättning moduler som går att komponera.

Alla tre skriver du själv: \code{reset_sync} och \code{button_sync} i Övning~\ref{c4:ex:resetsync}
och \ref{c4:ex:buttonsync}, \code{timer} i Övning~\ref{c7:ex:timer} i det här kapitlet. Kopiera in
dina egna i övningskatalogen innan du bygger den här. Att återanvända en modul förutsätter att man
har en, och det här är den första designen som ber dig om dem du redan byggt.

\begin{vhdlcode}
entity walking_led is
    generic(LED_COUNT : natural := 8;
            TICK_COUNT: natural := 25_000_000);
    port(clock, reset_n, button_n: in std_logic;
         led                     : out std_logic_vector(LED_COUNT-1 downto 0));
end entity;

architecture behaviour of walking_led is
signal reset_s2_n, button_edge_s2  : std_logic;
signal button_n_v, button_edge_s2_v: std_logic_vector(0 downto 0);
signal shift_enable, shift_tick    : std_logic;
signal shift_reg: std_logic_vector(LED_COUNT-1 downto 0);

begin
    led <= shift_reg;

    button_n_v(0)  <= button_n;
    button_edge_s2 <= button_edge_s2_v(0);

    reset_sync1: entity work.reset_sync
        port map(clock, reset_n, reset_s2_n);

    button_sync1: entity work.button_sync
        generic map(1)
        port map(clock, reset_s2_n, button_n_v, button_edge_s2_v);

    timer1: entity work.timer
        generic map(TICK_COUNT)
        port map(clock, reset_s2_n, shift_enable, shift_tick);

    SHIFT_ENABLE_PROCESS: process(clock, reset_s2_n) is
    begin
        if (reset_s2_n = '0') then
            shift_enable <= '0';
        elsif (rising_edge(clock)) then
            if (button_edge_s2 = '1') then
                shift_enable <= not shift_enable;
            end if;
        end if;
    end process;

    SHIFT_PROCESS: process(clock, reset_s2_n) is
    begin
        if (reset_s2_n = '0') then
            shift_reg <= (0 => '1', others => '0');
        elsif (rising_edge(clock)) then
            if (shift_tick = '1') then
                shift_reg <= shift_reg(LED_COUNT-2 downto 0) & shift_reg(LED_COUNT-1);
            end if;
        end if;
    end process;
end architecture;
\end{vhdlcode}

Designval värda att notera:
\begin{itemize}
  \item \code{button_n_v} och \code{button_edge_s2_v} överbryggar bara en breddskillnad:
    \code{walking_led}:s knappsignaler är enskilda bitar medan \code{button_sync}:s portar är
    vektorer med bredden \code{COUNT}. De två konkurrenta tilldelningarna kopplar skalären till
    element \code{0} och tillbaka. De kostar ingen hårdvara.
  \item Det här är ett PISO-register (\secref{c6:sec:piso}) i sin enklaste form: parallellt laddat
    exakt en gång, vid reset, med en enda tänd bit på position 0, och därefter bara skiftande. Det
    finns ingen \code{load} under körning, eftersom den här designen aldrig laddar ett nytt mönster.
  \item Det är ett \textbf{cirkulärt} skiftregister. Skiftningen är \secref{c6:sec:shift}:s uttryck
    oförändrat, så den tända biten vandrar mot MSB; den enda skillnaden mot ett vanligt PISO är vad
    som matar bit 0, vilket är biten som faller av toppen snarare än ett \code{serial_in} någon
    annanstans ifrån. Ett engångsregister som ”skiftar ut ett mönster” blir en oändligt upprepande
    mönstergenerator gratis, bara genom att återcirkulera sin egen utgång.
  \item \code{shift_enable} är \secref{c7:sec:compose}:s växlingsvippa återanvänd ordagrant:
    knappens synkroniserade flankpuls vänder en enda bit, och den biten blir timerns \code{enable}.
  \item \code{shift_tick} är timerns \code{timeout}-puls, återanvänd som skiftregistrets klockenable
    i stället för att driva en lysdiodsväxling. Samma byggblock, en annan konsument, och precis
    regeln ”en klocka, allt annat är enable” från \secref{c7:sec:compose}.
  \item Både \code{LED_COUNT} och \code{TICK_COUNT} är generics, så \code{generic map(4, 25_000_000)}
    ger en rad med 4 lysdioder och \code{generic map(8, 5_000_000)} ett snabbare steg på 100~ms,
    utan att modulen rörs.
\end{itemize}

På kortet, demonstrerat under passet, tilldelas portarna via Pin Planner (bilaga~\ref{app:quartus}):
\code{clock} till 50 MHz-oscillatorn, \code{reset_n} och \code{button_n} till två tryckknappar, och
\code{led} till en rad lysdioder på kortet. Tryck på knappen en gång för att sätta biten i rörelse,
en gång till för att frysa den.

Du kan se samma beteende utan kort genom att köra referenstestbänken, som skriver över båda generics
så att ett helt varv tar en handfull klockcykler:

\begin{shell}
cd lectures/L07/exercises/walking_led
cp ../../../L04/exercises/reset_sync/reset_sync.vhd .   # the three you wrote yourself
cp ../../../L04/exercises/button_sync/button_sync.vhd .
cp ../timer/timer.vhd .
ghdl -a --std=93 reset_sync.vhd button_sync.vhd timer.vhd \
                 walking_led.vhd walking_led_tb.vhd
ghdl -e --std=93 walking_led_tb
ghdl -r --std=93 walking_led_tb --assert-level=error --stop-time=10ms
\end{shell}
